\documentclass[a4paper,12pt,twoside]{memoir}

% Castellano
\usepackage[spanish,es-tabla]{babel}
\selectlanguage{spanish}
\usepackage[utf8]{inputenc}
\usepackage[T1]{fontenc}
\usepackage{lmodern} % scalable font
\usepackage{microtype}
\usepackage{placeins}

\RequirePackage{booktabs}
\RequirePackage[table]{xcolor}
\RequirePackage{xtab}
\RequirePackage{multirow}

% Links
\PassOptionsToPackage{hyphens}{url}\usepackage[colorlinks]{hyperref}
\hypersetup{
	allcolors = {red}
}

% Ecuaciones
\usepackage{amsmath}

% Rutas de fichero / paquete
\newcommand{\ruta}[1]{{\sffamily #1}}

% Párrafos
\nonzeroparskip

% Huérfanas y viudas
\widowpenalty100000
\clubpenalty100000

% Evitar solapes en el header
\nouppercaseheads

% Imagenes
\usepackage{graphicx}
\newcommand{\imagen}[2]{
	\begin{figure}[!h]
		\centering
		\includegraphics[width=0.9\textwidth]{#1}
		\caption{#2}\label{fig:#1}
	\end{figure}
	\FloatBarrier
}

\newcommand{\imagenflotante}[2]{
	\begin{figure}%[!h]
		\centering
		\includegraphics[width=0.9\textwidth]{#1}
		\caption{#2}\label{fig:#1}
	\end{figure}
}



% El comando \figura nos permite insertar figuras comodamente, y utilizando
% siempre el mismo formato. Los parametros son:
% 1 -> Porcentaje del ancho de página que ocupará la figura (de 0 a 1)
% 2 --> Fichero de la imagen
% 3 --> Texto a pie de imagen
% 4 --> Etiqueta (label) para referencias
% 5 --> Opciones que queramos pasarle al \includegraphics
% 6 --> Opciones de posicionamiento a pasarle a \begin{figure}
\newcommand{\figuraConPosicion}[6]{%
  \setlength{\anchoFloat}{#1\textwidth}%
  \addtolength{\anchoFloat}{-4\fboxsep}%
  \setlength{\anchoFigura}{\anchoFloat}%
  \begin{figure}[#6]
    \begin{center}%
      \Ovalbox{%
        \begin{minipage}{\anchoFloat}%
          \begin{center}%
            \includegraphics[width=\anchoFigura,#5]{#2}%
            \caption{#3}%
            \label{#4}%
          \end{center}%
        \end{minipage}
      }%
    \end{center}%
  \end{figure}%
}

%
% Comando para incluir imágenes en formato apaisado (sin marco).
\newcommand{\figuraApaisadaSinMarco}[5]{%
  \begin{figure}%
    \begin{center}%
    \includegraphics[angle=90,height=#1\textheight,#5]{#2}%
    \caption{#3}%
    \label{#4}%
    \end{center}%
  \end{figure}%
}
% Para las tablas
\newcommand{\otoprule}{\midrule [\heavyrulewidth]}
%
% Nuevo comando para tablas pequeñas (menos de una página).
\newcommand{\tablaSmall}[5]{%
 \begin{table}
  \begin{center}
   \rowcolors {2}{gray!35}{}
   \begin{tabular}{#2}
    \toprule
    #4
    \otoprule
    #5
    \bottomrule
   \end{tabular}
   \caption{#1}
   \label{tabla:#3}
  \end{center}
 \end{table}
}

%
%Para el float H de tablaSmallSinColores
\usepackage{float}

%
% Nuevo comando para tablas pequeñas (menos de una página).
\newcommand{\tablaSmallSinColores}[5]{%
 \begin{table}[H]
  \begin{center}
   \begin{tabular}{#2}
    \toprule
    #4
    \otoprule
    #5
    \bottomrule
   \end{tabular}
   \caption{#1}
   \label{tabla:#3}
  \end{center}
 \end{table}
}

\newcommand{\tablaApaisadaSmall}[5]{%
\begin{landscape}
  \begin{table}
   \begin{center}
    \rowcolors {2}{gray!35}{}
    \begin{tabular}{#2}
     \toprule
     #4
     \otoprule
     #5
     \bottomrule
    \end{tabular}
    \caption{#1}
    \label{tabla:#3}
   \end{center}
  \end{table}
\end{landscape}
}

%
% Nuevo comando para tablas grandes con cabecera y filas alternas coloreadas en gris.
\newcommand{\tabla}[6]{%
  \begin{center}
    \tablefirsthead{
      \toprule
      #5
      \otoprule
    }
    \tablehead{
      \multicolumn{#3}{l}{\small\sl continúa desde la página anterior}\\
      \toprule
      #5
      \otoprule
    }
    \tabletail{
      \hline
      \multicolumn{#3}{r}{\small\sl continúa en la página siguiente}\\
    }
    \tablelasttail{
      \hline
    }
    \bottomcaption{#1}
    \rowcolors {2}{gray!35}{}
    \begin{xtabular}{#2}
      #6
      \bottomrule
    \end{xtabular}
    \label{tabla:#4}
  \end{center}
}

%
% Nuevo comando para tablas grandes con cabecera.
\newcommand{\tablaSinColores}[6]{%
  \begin{center}
    \tablefirsthead{
      \toprule
      #5
      \otoprule
    }
    \tablehead{
      \multicolumn{#3}{l}{\small\sl continúa desde la página anterior}\\
      \toprule
      #5
      \otoprule
    }
    \tabletail{
      \hline
      \multicolumn{#3}{r}{\small\sl continúa en la página siguiente}\\
    }
    \tablelasttail{
      \hline
    }
    \bottomcaption{#1}
    \begin{xtabular}{#2}
      #6
      \bottomrule
    \end{xtabular}
    \label{tabla:#4}
  \end{center}
}

%
% Nuevo comando para tablas grandes sin cabecera.
\newcommand{\tablaSinCabecera}[5]{%
  \begin{center}
    \tablefirsthead{
      \toprule
    }
    \tablehead{
      \multicolumn{#3}{l}{\small\sl continúa desde la página anterior}\\
      \hline
    }
    \tabletail{
      \hline
      \multicolumn{#3}{r}{\small\sl continúa en la página siguiente}\\
    }
    \tablelasttail{
      \hline
    }
    \bottomcaption{#1}
  \begin{xtabular}{#2}
    #5
   \bottomrule
  \end{xtabular}
  \label{tabla:#4}
  \end{center}
}



\definecolor{cgoLight}{HTML}{EEEEEE}
\definecolor{cgoExtralight}{HTML}{FFFFFF}

%
% Nuevo comando para tablas grandes sin cabecera.
\newcommand{\tablaSinCabeceraConBandas}[5]{%
  \begin{center}
    \tablefirsthead{
      \toprule
    }
    \tablehead{
      \multicolumn{#3}{l}{\small\sl continúa desde la página anterior}\\
      \hline
    }
    \tabletail{
      \hline
      \multicolumn{#3}{r}{\small\sl continúa en la página siguiente}\\
    }
    \tablelasttail{
      \hline
    }
    \bottomcaption{#1}
    \rowcolors[]{1}{cgoExtralight}{cgoLight}

  \begin{xtabular}{#2}
    #5
   \bottomrule
  \end{xtabular}
  \label{tabla:#4}
  \end{center}
}




\graphicspath{ {./img/} }

% Capítulos
\chapterstyle{bianchi}
\newcommand{\capitulo}[2]{
	\setcounter{chapter}{#1}
	\setcounter{section}{0}
	\setcounter{figure}{0}
	\setcounter{table}{0}
	\chapter*{#2}
	\addcontentsline{toc}{chapter}{#2}
	\markboth{#2}{#2}
}

% Apéndices
\renewcommand{\appendixname}{Apéndice}
\renewcommand*\cftappendixname{\appendixname}

\newcommand{\apendice}[1]{
	%\renewcommand{\thechapter}{A}
	\chapter{#1}
}

\renewcommand*\cftappendixname{\appendixname\ }

% Formato de portada
\makeatletter
\usepackage{xcolor}
\newcommand{\tutor}[1]{\def\@tutor{#1}}
\newcommand{\course}[1]{\def\@course{#1}}
\definecolor{cpardoBox}{HTML}{E6E6FF}
\def\maketitle{
  \null
  \thispagestyle{empty}
  % Cabecera ----------------
\noindent\includegraphics[width=\textwidth]{cabecera}\vspace{1cm}%
  \vfill
  % Título proyecto y escudo informática ----------------
  \colorbox{cpardoBox}{%
    \begin{minipage}{.8\textwidth}
      \vspace{.5cm}\Large
      \begin{center}
      \textbf{TFG del Grado en Ingeniería Informática}\vspace{.6cm}\\
      \textbf{\LARGE\@title{}}
      \end{center}
      \vspace{.2cm}
    \end{minipage}

  }%
  \hfill\begin{minipage}{.20\textwidth}
    \includegraphics[width=\textwidth]{escudoInfor}
  \end{minipage}
  \vfill
  % Datos de alumno, curso y tutores ------------------
  \begin{center}%
  {%
    \noindent\LARGE
    Presentado por \@author{}\\ 
    en Universidad de Burgos --- \@date{}\\
    Tutor: \@tutor{}\\
  }%
  \end{center}%
  \null
  \cleardoublepage
  }
\makeatother


% Datos de portada
\title{título del TFG \\Documentación Técnica}
\author{nombre alumno}
\tutor{nombre tutor}
\date{\today}

\begin{document}

\maketitle



\cleardoublepage



%%%%%%%%%%%%%%%%%%%%%%%%%%%%%%%%%%%%%%%%%%%%%%%%%%%%%%%%%%%%%%%%%%%%%%%%%%%%%%%%%%%%%%%%



\frontmatter


\clearpage

% Indices
\tableofcontents

\clearpage

\listoffigures

\clearpage

\listoftables

\clearpage

\mainmatter

\appendix

\apendice{Plan de Proyecto Software}

\section{Introducción}

\section{Planificación temporal}

\section{Estudio de viabilidad}

\subsection{Viabilidad económica}

\subsection{Viabilidad legal}



\apendice{Especificación de Requisitos}

\section{Introducción}

En este apéndice se van a identificar y discutir los distintos objetivos que debe cumplir el proyecto, y los requisitos funcionales y no funcionales necesarios para que estos objetivos se lleguen a cumplir.


\section{Objetivos generales}

El objetivo del proyecto es el de disponer de una aplicación que permita mostrar a los consumidores de un centro de restauración los distintos menús que se ofrecen, que los usuarios puedan escoger el que deseen y que, dentro de él, a la hora de elegir los platos que van a consumir, puedan ver en tiempo real el cálculo de un informe nutricional en función de sus elecciones. 

Para escoger el menú que desean, los consumidores primero deben seleccionar una fecha para ver los menús disponibles en ese día.

Inicialmente el proyecto está planteado para ser usado por la Universidad de Burgos, pero se va a disponer de la opción de crear nuevas empresas, por lo que podría ser implementado en otros lugares. Cada empresa tiene una URL única, en la cual se muestran los distintos locales de esa empresa, por lo que esta URL podría ser la facilitada a los usuarios finales.

Además de esto, también se ofrece un panel de gestión en el que el personal de los centros pueda administrar los distintos recursos, y este panel ofrece distintas funciones de administración en función del rol que cumpla el personal en la empresa.

Existen tres roles diferenciados:

\begin{itemize}
	\item \textbf{Administrador}: Este tipo de usuario va a tener acceso para manejar los recursos globales de la aplicación, que son las empresas, los locales, los usuarios y los alimentos. Va a tener permisos para crear, editar o borrar cualquiera de estos 4 recursos. También va a tener posibilidad de ver los informes nutricionales de los alimentos. Es el único tipo de usuario que no va a estar vinculado a una empresa. 
	\item \textbf{Editor}: Este tipo de usuario va a tener acceso a los recursos vinculados a su empresa (exceptuando los locales), es decir, va a tener acceso para gestionar los menús y los platos, pudiendo crear nuevos elementos, editarlos o borrarlos. Además de esto, dentro del panel de creación de platos va a tener la posibilidad de crear nuevos alimentos, pero no va a poder editarlos ni borrarlos. También va a poder visualizar los informes nutricionales de los menús, los platos y los alimentos.
	\item \textbf{Camarero}: Este tipo de usuario sólo va a tener acceso a gestionar los menús, pudiendo crear, editar y borrar los distintos menús de su empresa. Va a poder visualizar los informes nutricionales de los menús.
\end{itemize}

Para el acceso al panel de gestión, existe una vista de inicio de sesión en la que el usuario debe indicar su nombre de usuario y contraseña. 

En caso de que el usuario haya olvidado su contraseña, hay un panel de contraseña olvidada en el que el usuario debe introducir su nombre de usuario, y a continuación, la nueva contraseña deseada.


\section{Catálogo de requisitos}

Se van a proceder a indicar los distintos requisitos, tanto funcionales como no funcionales:

\subsection{Requisitos funcionales}

\begin{itemize}
	\item \textbf{RF - 1}: Poder acceder al panel de gestión mediante el uso de credenciales
	\item \textbf{RF - 2}: Gestión de alimentos
	\begin{itemize}
		\item \textbf{RF - 2.1}: Búsqueda de alimento en Nutritionix
		\item \textbf{RF - 2.2}: Edición y eliminación de alimentos
		\begin{itemize}
			\item \textbf{RF - 2.2.1}: Edición de alimento
			\item \textbf{RF - 2.2.2}: Eliminación de alimento
		\end{itemize}
		\item \textbf{RF - 2.3}: Visualización de componentes nutricionales
		\begin{itemize}
			\item \textbf{RF - 2.3.1}: Gráfico
			\item \textbf{RF - 2.3.2}: Tabla nutricional
		\end{itemize}
		\item \textbf{RF - 2.4}: Creación de alimentos a mano 
	\end{itemize}
	
	\item \textbf{RF - 3}: Gestión de platos
	\begin{itemize}
		\item \textbf{RF - 3.1}: Creación, edición y eliminación de platos
		\begin{itemize}
			\item \textbf{RF - 3.1.1}: Creación de plato
			\item \textbf{RF - 3.1.2}: Edición de plato
			\item \textbf{RF - 3.1.3}: Eliminación de plato
		\end{itemize}
		\item \textbf{RF - 3.2}: Visualización de componentes nutricionales
		\begin{itemize}
			\item \textbf{RF - 3.2.1}: Gráfico
			\item \textbf{RF - 3.2.2}: Tabla nutricional
		\end{itemize}
		\item \textbf{RF - 3.3}: Visualización de ingredientes del plato 
	\end{itemize}
	
	\item \textbf{RF - 4}: Gestión de menús
	\begin{itemize}
		\item \textbf{RF - 4.1}: Creación, edición y eliminación de menús
		\begin{itemize}
			\item \textbf{RF - 4.1.1}: Creación de menús
			\item \textbf{RF - 4.1.2}: Edición de menús
			\item \textbf{RF - 4.1.3}: Eliminación de menús
		\end{itemize}
		\item \textbf{RF - 4.2}: Visualización de componentes nutricionales
		\begin{itemize}
			\item \textbf{RF - 4.2.1}: Gráfico
			\item \textbf{RF - 4.2.2}: Tabla nutricional
			\item \textbf{RF - 4.2.3}: Selección de platos del menú
		\end{itemize}
		\item \textbf{RF - 4.3}: Visualización de platos del menú
	\end{itemize}
	
	\item \textbf{RF - 5}: Gestión de usuarios
	\begin{itemize}
		\item \textbf{RF - 5.1}: Creación de usuarios
		\item \textbf{RF - 5.2}: Edición de usuarios
		\item \textbf{RF - 5.3}: Eliminación de usuarios
	\end{itemize}
	
	\item \textbf{RF - 6}: Gestión de empresas
	\begin{itemize}
		\item \textbf{RF - 6.1}: Creación de empresas
		\item \textbf{RF - 6.2}: Edición de empresas
		\item \textbf{RF - 6.3}: Eliminación de empresas
	\end{itemize}
	
	\item \textbf{RF - 7}: Gestión de locales
	\begin{itemize}
		\item \textbf{RF - 7.1}: Creación de locales
		\item \textbf{RF - 7.2}: Edición de locales
		\item \textbf{RF - 7.3}: Eliminación de locales
	\end{itemize}
	
	\item \textbf{RF - 8}: Restablecimiento de contraseña
	\item \textbf{RF - 9}: Muestra de contenido al usuario final
	\begin{itemize}
		\item \textbf{RF - 9.1}: Listado de empresas
		\item \textbf{RF - 9.2}: Listado de locales
		\item \textbf{RF - 9.3}: Selección de fecha para ver menús publicados
		\item \textbf{RF - 9.4}: Listado de menús en una fecha
	\end{itemize}
\end{itemize}

\subsection{Requisitos no funcionales}

\begin{itemize}
\item \textbf{RNF - 1. Eficiencia}: Se debe tener una respuesta rápida a las peticiones del usuario, tanto por parte de la propia aplicación como por parte de la API REST, que debe manejar la información en el menor tiempo posible.
\item \textbf{RNF - 2. Escalabilidad}: La aplicación debe ser capaz de soportar grandes cargas de usuarios, puesto que al final la mayoría de usuarios van a usarla en la misma franja de tiempo, la hora de la comida.
\item \textbf{RNF - 3. Modularidad}: La aplicación debe estar formada por módulos independientes, que se puedan reemplazar sin afectar al resto de la funcionalidad de la aplicación.
\item \textbf{RNF - 4. Usabilidad}: Debe ser un producto fácil de usar, que no requiera de conocimientos avanzados de informática. Las distintas funciones deben ser intuitivas y estar claramente definidas, sin que haga falta indicar cómo realizarlas de forma explícita.
\item \textbf{RNF - 5. Responsibilidad}: La aplicación debe ser adaptable a los distintos tamaños de pantalla, funcionando igual de bien tanto en ordenadores como en dispositivos móviles.
\item \textbf{RNF - 6. Legalidad}: Se deben cumplir con las normativas y leyes del territorio en el que se ejecute la aplicación, que de momento sólo es España.
\end{itemize}


\section{Especificación de requisitos}

\begin{table}[p]
	\centering
	\begin{tabularx}{\linewidth}{ p{0.21\columnwidth} p{0.71\columnwidth} }
		\toprule
		\textbf{CU-1}    & \textbf{Ejemplo de caso de uso}\\
		\toprule
		\textbf{Versión}              & 1.0    \\
		\textbf{Autor}                & Alumno \\
		\textbf{Requisitos asociados} & RF-xx, RF-xx \\
		\textbf{Descripción}          & La descripción del CU \\
		\textbf{Precondición}         & Precondiciones (podría haber más de una) \\
		\textbf{Acciones}             &
		\begin{enumerate}
			\def\labelenumi{\arabic{enumi}.}
			\tightlist
			\item Pasos del CU
			\item Pasos del CU (añadir tantos como sean necesarios)
		\end{enumerate}\\
		\textbf{Postcondición}        & Postcondiciones (podría haber más de una) \\
		\textbf{Excepciones}          & Excepciones \\
		\textbf{Importancia}          & Alta o Media o Baja... \\
		\bottomrule
	\end{tabularx}
	\caption{CU-1 Nombre del caso de uso.}
\end{table}

% TODO Especificación de requisitos
\apendice{Especificación de diseño}

\section{Introducción}

\section{Diseño de datos}

\section{Diseño procedimental}

Uno de los grandes cambios que se han realizado en este proyecto respecto a la iteración anterior es el desarrollo de una API REST completa para trabajar e interactuar con los datos. A continuación voy a proceder a detallar y documentar cada uno de los distintos endpoints a los cuáles se puede hacer llamadas HTTP:

% TODO Indicar cómo hacer llamadas a la API %
% TODO Indicar el endpoint al que hacer las llamadas %

\subsection{Empresas}

\begin{apiRoute}{get}{/empresas}{Devuelve todas las empresas existentes}
	
	\begin{routeParameter}
		\routeParamItem{id\_empresa}{Devuelve la empresa con el ID correspondiente}
		\routeParamItem{nombre}{Devuelve todas las empresas cuyo nombre contengan el valor pasado por parámetro}
		\routeParamItem{email}{Devuelve la empresa con el email correspondiente}
		\routeParamItem{telefono}{Devuelve la empresa con el teléfono correspondiente}
		\routeParamItem{cif}{Devuelve la empresa con el CIF correspondiente}
	\end{routeParameter}
	
	\begin{routeResponse}{application/json}
		\begin{routeResponseItem}{200}{ok}
			\begin{routeResponseItemBody}
{
	"id": 1,
	"nombre": "UBU",
	"email": "cafeterias@ubu.es",
	"direccion": "Calle Don Juan de Austria 1, 09001 Burgos",
	"telefono": "947258060",
	"cif": "123456789"
}
			\end{routeResponseItemBody}
		\end{routeResponseItem}
		\begin{routeResponseItem}{400}{Bad Request: EntityDoesntExistsException}
			\begin{routeResponseItemBody}
{
    "timestamp": "2023-09-19T17:18:27.228+00:00",
    "status": 400,
    "error": "Bad Request",
    "trace": "ooo.alvar.nutrimenu.apirest.excepciones...",
    "message": "No existe una empresa con id 3",
    "path": "/empresas"
}
			\end{routeResponseItemBody}
		\end{routeResponseItem}
	\end{routeResponse}
	
\end{apiRoute}

\begin{apiRoute}{post}{/empresas}{Añade una nueva empresa}
	\begin{routeParameter}
		\noRouteParameter{No hay parámetros}
	\end{routeParameter}
	\begin{routeRequest}{application/json}
		\begin{routeRequestBody}
{
	"nombre": "Vips",
	"email": "empresa@vips.com",
	"direccion": "Puerta del Sol, Madrid",
	"telefono": "900123123",
	"cif": "123456787"
}
		\end{routeRequestBody}
	\end{routeRequest}
	\begin{routeResponse}{application/json}
		\begin{routeResponseItem}{200}{ok}
			\begin{routeResponseItemBody}
{
	"id": 3,
	"nombre": "Vips",
	"email": "empresa@vips.com",
	"direccion": "Puerta del Sol, Madrid",
	"telefono": "900123123",
	"cif": "123456787"
}
			\end{routeResponseItemBody}
		\end{routeResponseItem}
	\end{routeResponse}
\end{apiRoute}

\begin{apiRoute}{put}{/empresas}{Actualiza la información de una empresa ya existente}
	\begin{routeParameter}
		\routeParamItem{id\_empresa}{Modifica cualquier valor de la empresa con el ID correspondiente}
	\end{routeParameter}
	
	\begin{routeRequest}{application/json}
		\begin{routeRequestBody}
{
	"nombre": "Universidad de Burgos"
}
		\end{routeRequestBody}
	\end{routeRequest}
	\begin{routeResponse}{application/json}
		\begin{routeResponseItem}{200}{ok}
			\begin{routeResponseItemBody}
{
	"id": 1,
	"nombre": "UBU",
	"email": "cafeterias@ubu.es",
	"direccion": "Calle Don Juan de Austria 1, 09001 Burgos",
	"telefono": "947258060",
	"cif": "123456789"
}
			\end{routeResponseItemBody}
		\end{routeResponseItem}
		\begin{routeResponseItem}{400}{Bad Request: EntityDoesntExistsException}
			\begin{routeResponseItemBody}
{
    "timestamp": "2023-09-19T17:18:27.228+00:00",
    "status": 400,
    "error": "Bad Request",
    "trace": "ooo.alvar.nutrimenu.apirest.excepciones...",
    "message": "No existe una empresa con id 3",
    "path": "/empresas"
}
			\end{routeResponseItemBody}
		\end{routeResponseItem}
	\end{routeResponse}
\end{apiRoute}

\begin{apiRoute}{delete}{/empresas}{Elimina una empresa de la base de datos}
	\begin{routeParameter}
		\routeParamItem{id\_empresa}{Elimina la empresa con el ID correspondiente}
	\end{routeParameter}
	\begin{routeResponse}{application/json}
		\begin{routeResponseItem}{200}{ok}
			\begin{routeResponseItemBody}
Empresa con id eliminada correctamente	
			\end{routeResponseItemBody}
		\end{routeResponseItem}
		
\begin{routeResponseItem}{400}{Bad Request: EntityDoesntExistsException}
			\begin{routeResponseItemBody}
{
    "timestamp": "2023-09-19T17:18:27.228+00:00",
    "status": 400,
    "error": "Bad Request",
    "trace": "ooo.alvar.nutrimenu.apirest.excepciones...",
    "message": "No existe una empresa con id 3",
    "path": "/empresas"
}
			\end{routeResponseItemBody}
		\end{routeResponseItem}
		
	\end{routeResponse}
\end{apiRoute}

\subsection{Locales}

\begin{apiRoute}{get}{/locales}{Devuelve todos los locales existentes}
	
	\begin{routeParameter}
		\routeParamItem{id\_local}{Devuelve el local con el ID correspondiente}
		\routeParamItem{id\_empresa}{Devuelve todos los locales correspondientes a la empresa con el ID correspondiente}
		\routeParamItem{nombre}{Devuelve todos los locales cuyo nombre contengan el valor pasado por parámetro}
		\routeParamItem{email}{Devuelve el local con el email correspondiente}
		\routeParamItem{telefono}{Devuelve el local con el teléfono correspondiente}
	\end{routeParameter}
	
	\begin{routeResponse}{application/json}
		\begin{routeResponseItem}{200}{ok}
			\begin{routeResponseItemBody}
{
	"id": 1,
	"nombre": "Escuela Politecnica",
	"email": "eps@ubu.es",
	"direccion": "Avenida Cantabria",
	"telefono": "947123123",
	"empresa": {
		"id": 1,
		"nombre": "Universidad de Burgos",
		"email": "cafeterias@ubu.es",
		"direccion": "Calle Don Juan de Austria 1, 09001 Burgos",
		"telefono": "9002375432",
		"cif": "1113111"
        }
    }
			\end{routeResponseItemBody}
		\end{routeResponseItem}
		\begin{routeResponseItem}{400}{Bad Request: EntityDoesntExistsException}
			\begin{routeResponseItemBody}
{
    "timestamp": "2023-09-19T17:18:27.228+00:00",
    "status": 400,
    "error": "Bad Request",
    "trace": "ooo.alvar.nutrimenu.apirest.excepciones...",
    "message": "No existe una empresa con id 3",
    "path": "/locales"
}
			\end{routeResponseItemBody}
		\end{routeResponseItem}
	\end{routeResponse}
	
\end{apiRoute}

\begin{apiRoute}{post}{/locales}{Añade un nuevo local}
	\begin{routeParameter}
		\routeParamItem{id\_empresa}{Indica a qué empresa pertenece el local}
	\end{routeParameter}
	\begin{routeRequest}{application/json}
		\begin{routeRequestBody}
{
	"nombre": "Cafeteria EPS",
	"email": "eps@ubu.es",
	"direccion": "Avenida Cantabria",
	"telefono": "947123954"
}
		\end{routeRequestBody}
	\end{routeRequest}
	\begin{routeResponse}{application/json}
		\begin{routeResponseItem}{200}{ok}
			\begin{routeResponseItemBody}
{
    "id": 3,
    "nombre": "Cafeteria EPS",
    "email": "eps@ubu.es",
    "direccion": "Avenida Cantabria",
    "telefono": "947123954",
    "empresa": {
        "id": 1,
        "nombre": "Universidad de Burgos",
        "email": "cafeterias@ubu.es",
        "direccion": "Calle Don Juan de Austria 1, 09001 Burgos",
        "telefono": "9002375432",
        "cif": "1113111"
    }
}
			\end{routeResponseItemBody}
		\end{routeResponseItem}
		\begin{routeResponseItem}{400}{Bad Request: EntityDoesntExistsException}
			\begin{routeResponseItemBody}
{
    "timestamp": "2023-09-19T17:18:27.228+00:00",
    "status": 400,
    "error": "Bad Request",
    "trace": "ooo.alvar.nutrimenu.apirest.excepciones...",
    "message": "No existe una empresa con id 3",
    "path": "/locales"
}
			\end{routeResponseItemBody}
		\end{routeResponseItem}
	\end{routeResponse}
\end{apiRoute}

\begin{apiRoute}{put}{/locales}{Actualiza la información de un local ya existente}
	\begin{routeParameter}
		\routeParamItem{id\_local}{Modifica cualquier valor del local con el ID correspondiente}
	\end{routeParameter}
	
	\begin{routeRequest}{application/json}
		\begin{routeRequestBody}
{
	"nombre": "Cafeteria Economicas"
}
		\end{routeRequestBody}
	\end{routeRequest}
	\begin{routeResponse}{application/json}
		\begin{routeResponseItem}{200}{ok}
			\begin{routeResponseItemBody}
{
    "id": 3,
    "nombre": "Cafeteria Economicas",
    "email": "eps@ubu.es",
    "direccion": "Avenida Cantabria",
    "telefono": "947123954",
    "empresa": {
        "id": 1,
        "nombre": "Universidad de Burgos",
        "email": "cafeterias@ubu.es",
        "direccion": "Calle Don Juan de Austria 1, 09001 Burgos",
        "telefono": "9002375432",
        "cif": "1113111"
    }
}
			\end{routeResponseItemBody}
		\end{routeResponseItem}
		\begin{routeResponseItem}{400}{Bad Request: EntityDoesntExistsException}
			\begin{routeResponseItemBody}
{
    "timestamp": "2023-09-19T17:18:27.228+00:00",
    "status": 400,
    "error": "Bad Request",
    "trace": "ooo.alvar.nutrimenu.apirest.excepciones...",
    "message": "No existe un local con id 3",
    "path": "/locales"
}
			\end{routeResponseItemBody}
		\end{routeResponseItem}
	\end{routeResponse}
\end{apiRoute}

\begin{apiRoute}{put}{/add/local/menus}{Añade un menú a un local}
	\begin{routeParameter}
		\routeParamItem{id\_local}{ID del local al que se le va a añadir un menú}
		\routeParamItem{id\_menu}{ID del menú que va a ser añadido al local}
	\end{routeParameter}
	
	\begin{routeResponse}{application/json}
		\begin{routeResponseItem}{200}{ok}
			\begin{routeResponseItemBody}
{
    "id": 1,
    "nombre": "Menu del miercoles",
    "descripcion": "Menu del martes",
    "fechaCreacion": "2023-09-19T19:00:14.998139Z",
    "fechaModificacion": "2023-09-19T19:22:57.314517Z",
    "platos": [
        {
            "id": 1,
            "tipoPlato": "PRIMER_PLATO",
            "nombre": "Kevin Bacon",
            "descripcion": "Hamburguesa rellena de bacon",
            "fechaCreacion": "2023-09-19T19:00:59.387553Z",
            "fechaModificacion": "2023-09-19T19:00:59.387556Z"
        }
    ]
}
			\end{routeResponseItemBody}
		\end{routeResponseItem}
		\begin{routeResponseItem}{400}{Bad Request: EntityDoesntExistsException}
			\begin{routeResponseItemBody}
{
    "timestamp": "2023-09-19T17:18:27.228+00:00",
    "status": 400,
    "error": "Bad Request",
    "trace": "ooo.alvar.nutrimenu.apirest.excepciones...",
    "message": "No existe un local con id 3",
    "path": "/add/local/menus"
}
			\end{routeResponseItemBody}
		\end{routeResponseItem}
	\end{routeResponse}
\end{apiRoute}

\begin{apiRoute}{delete}{/locales}{Elimina un local de la base de datos}
	\begin{routeParameter}
		\routeParamItem{id\_local}{Elimina el local con el ID correspondiente}
	\end{routeParameter}
	\begin{routeResponse}{application/json}
		\begin{routeResponseItem}{200}{ok}
			\begin{routeResponseItemBody}
Local con id eliminado correctamente	
			\end{routeResponseItemBody}
		\end{routeResponseItem}
		
\begin{routeResponseItem}{400}{Bad Request: EntityDoesntExistsException}
			\begin{routeResponseItemBody}
{
    "timestamp": "2023-09-19T17:18:27.228+00:00",
    "status": 400,
    "error": "Bad Request",
    "trace": "ooo.alvar.nutrimenu.apirest.excepciones...",
    "message": "No existe un local con id 3",
    "path": "/locales"
}
			\end{routeResponseItemBody}
		\end{routeResponseItem}
		
	\end{routeResponse}
\end{apiRoute}

\subsection{Menús}

\begin{apiRoute}{get}{/menus}{Devuelve todos los menús existentes}
	
	\begin{routeParameter}
		\routeParamItem{id\_menu}{Devuelve el menú con el ID correspondiente}
		\routeParamItem{id\_local}{Devuelve todos los menús del local con el ID correspondiente}
	\end{routeParameter}
	
	\begin{routeResponse}{application/json}
		\begin{routeResponseItem}{200}{ok}
			\begin{routeResponseItemBody}
{
	"id": 1,
	"nombre": "Menu del martes",
	"descripcion": "Menu del martes",
	"fechaCreacion": "2023-09-19T19:00:14.998139Z",
	"fechaModificacion": "2023-09-19T19:00:14.998140Z",
	"platos": []
}
			\end{routeResponseItemBody}
		\end{routeResponseItem}
		\begin{routeResponseItem}{400}{Bad Request: EntityDoesntExistsException}
			\begin{routeResponseItemBody}
{
    "timestamp": "2023-09-19T17:18:27.228+00:00",
    "status": 400,
    "error": "Bad Request",
    "trace": "ooo.alvar.nutrimenu.apirest.excepciones...",
    "message": "No existe un menu con id 3",
    "path": "/menus"
}
			\end{routeResponseItemBody}
		\end{routeResponseItem}
		\begin{routeResponseItem}{400}{Bad Request: LackOfParametersException}
			\begin{routeResponseItemBody}
No se ha especificado ningun parametro de busqueda
			\end{routeResponseItemBody}
		\end{routeResponseItem}
	\end{routeResponse}
	
\end{apiRoute}

\begin{apiRoute}{post}{/menus}{Añade un nuevo menú}
	\begin{routeParameter}
		\routeParamItem{id\_local}{Indica a qué local pertenece el menú}
	\end{routeParameter}
	\begin{routeRequest}{application/json}
		\begin{routeRequestBody}
{
	"nombre": "Menu del martes",
	"descripcion": "Menu del martes"
}
		\end{routeRequestBody}
	\end{routeRequest}
	\begin{routeResponse}{application/json}
		\begin{routeResponseItem}{200}{ok}
			\begin{routeResponseItemBody}
{
	"id": 1,
	"nombre": "Menu del martes",
	"descripcion": "Menu del martes",
	"fechaCreacion": "2023-09-19T19:00:14.998139Z",
	"fechaModificacion": "2023-09-19T19:00:14.998140Z",
	"platos": []
}
			\end{routeResponseItemBody}
		\end{routeResponseItem}
		\begin{routeResponseItem}{400}{Bad Request: EntityDoesntExistsException}
			\begin{routeResponseItemBody}
{
    "timestamp": "2023-09-19T17:18:27.228+00:00",
    "status": 400,
    "error": "Bad Request",
    "trace": "ooo.alvar.nutrimenu.apirest.excepciones...",
    "message": "No existe un local con id 3",
    "path": "/menus"
}
			\end{routeResponseItemBody}
		\end{routeResponseItem}
	\end{routeResponse}
\end{apiRoute}

\begin{apiRoute}{put}{/menus}{Actualiza la información de un menú ya existente}
	\begin{routeParameter}
		\routeParamItem{id\_empresa}{Modifica cualquier valor del menú con el ID correspondiente}
	\end{routeParameter}
	
	\begin{routeRequest}{application/json}
		\begin{routeRequestBody}
{
	"nombre": "Menu del miercoles"
}
		\end{routeRequestBody}
	\end{routeRequest}
	\begin{routeResponse}{application/json}
		\begin{routeResponseItem}{200}{ok}
			\begin{routeResponseItemBody}
{
	"id": 1,
	"nombre": "Menu del miercoles",
	"descripcion": "Menu del martes",
	"fechaCreacion": "2023-09-19T19:00:14.998139Z",
	"fechaModificacion": "2023-09-19T19:00:14.998140Z",
	"platos": []
}
			\end{routeResponseItemBody}
		\end{routeResponseItem}
		\begin{routeResponseItem}{400}{Bad Request: EntityDoesntExistsException}
			\begin{routeResponseItemBody}
{
    "timestamp": "2023-09-19T17:18:27.228+00:00",
    "status": 400,
    "error": "Bad Request",
    "trace": "ooo.alvar.nutrimenu.apirest.excepciones...",
    "message": "No existe un menu con id 3",
    "path": "/menus"
}
			\end{routeResponseItemBody}
		\end{routeResponseItem}
	\end{routeResponse}
\end{apiRoute}

\begin{apiRoute}{put}{/add/menu/platos}{Añade un plato a un menú}
	\begin{routeParameter}
		\routeParamItem{id\_menu}{ID del menú al que se le va a añadir un plato}
		\routeParamItem{id\_plato}{ID del plato que va a ser añadido al menú}
	\end{routeParameter}
	
	\begin{routeResponse}{application/json}
		\begin{routeResponseItem}{200}{ok}
			\begin{routeResponseItemBody}
{
    "id": 1,
    "nombre": "Menu del miercoles",
    "descripcion": "Menu del martes",
    "fechaCreacion": "2023-09-19T19:00:14.998139Z",
    "fechaModificacion": "2023-09-19T19:22:57.314517Z",
    "platos": [
        {
            "id": 1,
            "tipoPlato": "PRIMER_PLATO",
            "nombre": "Kevin Bacon",
            "descripcion": "Hamburguesa rellena de bacon",
            "fechaCreacion": "2023-09-19T19:00:59.387553Z",
            "fechaModificacion": "2023-09-19T19:00:59.387556Z"
        }
    ]
}
			\end{routeResponseItemBody}
		\end{routeResponseItem}
		\begin{routeResponseItem}{400}{Bad Request: EntityDoesntExistsException}
			\begin{routeResponseItemBody}
{
    "timestamp": "2023-09-19T17:18:27.228+00:00",
    "status": 400,
    "error": "Bad Request",
    "trace": "ooo.alvar.nutrimenu.apirest.excepciones...",
    "message": "No existe un menu con id 3",
    "path": "/add/menu/platos"
}
			\end{routeResponseItemBody}
		\end{routeResponseItem}
	\end{routeResponse}
\end{apiRoute}

\begin{apiRoute}{delete}{/menus}{Elimina un menú de la base de datos}
	\begin{routeParameter}
		\routeParamItem{id\_menu}{Elimina el menú con el ID correspondiente}
	\end{routeParameter}
	\begin{routeResponse}{application/json}
		\begin{routeResponseItem}{200}{ok}
			\begin{routeResponseItemBody}
Menu con id eliminado correctamente	
			\end{routeResponseItemBody}
		\end{routeResponseItem}
		
\begin{routeResponseItem}{400}{Bad Request: EntityDoesntExistsException}
			\begin{routeResponseItemBody}
{
    "timestamp": "2023-09-19T17:18:27.228+00:00",
    "status": 400,
    "error": "Bad Request",
    "trace": "ooo.alvar.nutrimenu.apirest.excepciones...",
    "message": "No existe un menu con id 3",
    "path": "/menus"
}
			\end{routeResponseItemBody}
		\end{routeResponseItem}
		
	\end{routeResponse}
\end{apiRoute}

\section{Diseño arquitectónico}



\apendice{Documentación técnica de programación}

\section{Introducción}

\section{Estructura de directorios}

El proyecto se encuentra alojado en un repositorio de GitHub, el cuál contiene tanto esta documentación, como las dos aplicaciones principales (API REST y frontend), así como toda la infraestructura necesaria para ser desplegadas. 

La URL del repositorio es la siguiente:

\href{https://github.com/alvaromanjon/NutriMenu}{https://github.com/alvaromanjon/NutriMenu}

\imagen{GitHub/clone-repo}{Descarga del repositorio del proyecto}{0.8}

Podemos descargar el repositorio desde el propio GitHub, haciendo click en el botón \textit{Code} > \textit{Download ZIP}, o clonando el repositorio escribiendo en una línea de comandos:

\begin{lstlisting}[language=Bash]
  git clone https://github.com/alvaromanjon/NutriMenu.git
\end{lstlisting}

\subsection{Estructura de los archivos de Docker}

\dirtree{%
.1 NutriMenu.
.2 data.
.3 nutridb.sql.
.2 docker-compose.
.3 .env.
.3 .env.example.
.3 docker-compose.ci.yml.
.3 docker-compose.dev.yml.
.2 webapps.
.3 apirest.
.4 Dockerfile.
.4 run\_docker\_ci.sh.
.4 run\_docker\_dev.sh.
.3 frontend.
.4 .dockerignore.
.4 Dockerfile.
}

\section{Manual del programador}

\subsection{Dockerfiles}

\subsubsection{API REST}

\begin{lstlisting}
FROM eclipse-temurin:17
RUN apt update && apt -y upgrade
RUN apt install -y inotify-tools dos2unix
ENV HOME=/app
RUN mkdir -p $HOME
WORKDIR $HOME
\end{lstlisting}

A la hora de desplegar la parte de \textit{backend} disponemos para ello de un \textbf{Dockerfile}, que va a ser el encargado de definir cómo se va a crear la imagen de Docker, y qué parámetros y dependencias va a necesitar. Es el equivalente al proceso al que todos estamos acostumbrados de crear una máquina virtual, instalar el sistema operativo escogido y configurar todas las dependencias y paquetes necesarios, solo que en este caso todo esto podemos ahorrárnoslo (tanto en esfuerzo como en tiempo, ya que el tiempo de despliegue de una imagen de Docker es de apenas un par de minutos en la mayoría de los casos) escribiendo en este archivo la configuración necesaria que necesitamos para nuestro caso de uso.

En este caso me he basado en otra imagen ya existente, \textit{\textbf{eclipse-temurin}} (\href{https://hub.docker.com/layers/library/eclipse-temurin/17/images/sha256-b0faf02bf7acfc65be1c2d0a291140300bd129620f145bf1013a1da748295d0c?context=explore}{Docker Hub}), ya que contiene tanto el JDK como el JRE de Java, además de una instalación mínima de \textbf{Linux Ubuntu} con todas las dependencias ya instaladas y configuradas. He especificado que quiero usar la etiqueta \textit{\textbf{17}}, puesto que quiero usar Java 17 para este proyecto, ya que es la última versión LTS \textit{(Long Term Support)} disponible en el momento \cite{java:releases}. Existen muchas implementaciones del JDK de Java \cite{foojay:java-jdks}, pero en este caso he usado Eclipse Temurin principalmente por su soporte a arquitecturas \textit{ARMv7} y \textit{ARM64/v8}, lo que permite poder ejecutar los contenedores en plataformas como Apple Silicon, o incluso en mini-computadores como una Raspberry Pi.

Aunque la imagen descarga automáticamente la última versión disponible de \href{https://hub.docker.com/layers/library/eclipse-temurin/17/images/sha256-b0faf02bf7acfc65be1c2d0a291140300bd129620f145bf1013a1da748295d0c?context=explore}{Docker Hub}, por si acaso cada vez que construyamos la imagen vamos a comprobar si hay actualizaciones disponibles para cualquier paquete del sistema operativo del que dispone la imagen. Además de esto, también se van a instalar como dependencias \textit{\href{https://github.com/inotify-tools}{inotify-tools}}, una herramienta para monitorizar cambios en el código; y \textit{\href{https://dos2unix.sourceforge.io/}{dos2unix}}, que se encarga de cambiar el tipo de final de línea de \textbf{CRLF} (Windows) a \textbf{LF} (Unix) \cite{medium:spring-boot-in-docker}. Indicamos también que, dentro de la imagen, el directorio que vamos a usar para trabajar es \verb,/app,, puesto que esto nos permite que cualquier comando que ejecutemos a continuación se realice dentro de ese directorio \cite{docker:workdir}.

\subsubsection{Frontend} \label{dockerfile-frontend}

\begin{lstlisting}
FROM node:alpine
ENV HOME=/app
RUN mkdir -p $HOME
WORKDIR $HOME
COPY package*.json ./
RUN npm install
EXPOSE 3000
CMD ["npm", "start"]
\end{lstlisting}

En la parte del \textit{frontend} también disponemos de otro \textbf{Dockerfile}, en este caso con la imagen oficial de \textit{\textbf{Node}}, ya que es el entorno de ejecución que hay debajo de React. He usado la etiqueta \textit{\textbf{alpine}} (\href{https://hub.docker.com/layers/library/node/alpine/images/sha256-923cd6fac65f6892aa8bbb4208ad708c56b35f9ab86eca07ccc7b56dd28c9c77?context=explore}{Docker Hub}), ya que \textbf{Alpine} es una distribución de Linux orientada a ser lo más ligera posible, reduciendo el número de paquetes al mínimo y sustituyendo las herramientas GNU por \textbf{BusyBox} \cite{wikipedia:busybox}, un ejecutable que es capaz de emularlas.

Al igual que en el otro Dockerfile que acabamos de ver, he configurado el entorno de ejecución en el directorio \verb,/app, del contenedor. A continuación he indicado que quiero que se copien los archivos que empiezan por \verb,package, y acaban en \verb,.json, al contenedor, ya que en un proyecto React son los archivos que indican todas las dependencias que necesita el proyecto para poder ejecutarse. Sobre el qué directorio va a usar Docker para coger estos archivos, lo veremos después en el archivo de configuración de \textbf{Docker Compose}, puesto que al construir la imagen con la etiqueta \verb,build:, indicaremos el directorio al que puede acceder Docker durante la construcción de la imagen con el atributo \verb,context:, \cite{docker:build-context}.

Por último, la imagen va a instalar todas las dependencias del proyecto necesarias con \verb,npm install,, se va a abrir el puerto \textit{3000} del contenedor, ya que es el utilizado por React para poder acceder a la aplicación, y se va a iniciar la aplicación con el comando \verb,npm start,. 

\subsection{Docker Compose}

Dentro de los archivos de configuración de Docker Compose, situados en la carpeta \verb,docker-compose, del directorio raíz del proyecto, tenemos dos archivos distintos:

\begin{itemize}
	\item \verb,docker-compose.ci.yml,: Es el encargado de configurar los contenedores que se van a desplegar en el \textbf{entorno CI/CD configurado}, es decir, \textit{GitHub Actions}. Es prácticamente idéntico al entorno de desarrollo, tan sólo cambian las referencias a los scripts de arranque, como veremos a continuación.
	\item \verb,docker-compose.dev.yml,: Es el \textbf{entorno de desarrollo} utilizado durante la creación de esta aplicación, ya que está configurado de tal forma que permite interactuar con la aplicación y aplicar y compilar los cambios en el código en tiempo real, sin necesidad de reiniciar los contenedores, usando \textit{LiveReload} \cite{codejava:livereload}.
\end{itemize}

\subsubsection{Frontend}

\begin{lstlisting}
  frontend:
    build:
      context: ../webapps/frontend
    restart: always
    ports:
      - ${FRONTEND_DOCKER_PORT}:3000
    networks:
      - nutrimenu-net
    env_file:
      - .env
    depends_on:
      - api-rest
    volumes:
      - ../webapps/frontend:/app:rw
      - /app/node_modules
    stdin_open: true	
\end{lstlisting}

Dentro de nuestro archivo de configuración \verb,docker-compose.dev.yml, el primer servicio con el que nos encontramos es el correspondiente al contenedor encargado de desplegar la parte del \textit{frontend}. Esta parte no aparece dentro de \verb,docker-compose.ci.yml,, ya que GitHub Actions sólo está configurado para la parte del \textit{backend} y la base de datos. 

Dentro de este archivo de configuración podemos destacar algunos elementos:

\begin{itemize}
	\item Como hemos mencionado anteriormente en el apartado de los \textit{Dockerfile} \ref{dockerfile-frontend}, la etiqueta \verb,context:, en \verb,build:, indica el \textbf{directorio al que vamos a dar acceso a Docker a la hora de construir la imagen}, que en este caso es \verb,webapps/frontend,. Los dos puntos antes del directorio significan que ese directorio se encuentra fuera de la carpeta actual, \verb,docker-compose,.
	\item El contenedor está configurado para \textbf{reiniciarse de forma automática} en caso de fallo \cite{baeldung:restart-policies}.
	\item El contenedor \textbf{va a estar conectado a una red} creada también en el fichero de Docker Compose, la cual veremos a continuación.
	\item Se va a usar el fichero de \textbf{variables de entorno} descrito en el apartado \ref{entorno-desarrollo} para definir las distintas variables especificadas en este archivo.
	\item El contenedor va a esperar a que el contenedor de la \textbf{API REST} esté en estado \textit{ready} para arrancar \cite{docker:depends-on}.
	\item \textbf{Se van a crear dos volúmenes}, uno encargado de \textbf{mapear los archivos de la aplicación en local} (\verb,../webapps/frontend,) \textbf{con el directorio de trabajo dentro del contenedor} (\verb,/app,) ,indicando que este directorio va a ser tanto de lectura como de escritura; y otro volumen específico para \verb,node_modules,, \textbf{la carpeta que contiene todas las dependencias descargadas}. Este primer volumen es lo que nos permite poder realizar modificaciones sin tener que reiniciar servicios ni volver a construir imágenes, mientras que el segundo volumen se monta debido a que, cuando montamos este primer volumen, el contenido que hubiera en nuestro directorio de trabajo del contenedor (\verb,/app,) se sobreescribe por el del directorio mapeado, así que con este segundo volumen montado podemos recuperar la carpeta \verb,node_modules, generada al crear la imagen con el comando \verb,npm install,, puesto que sin esta carpeta \textbf{no podemos ejecutar la aplicación} (a no ser que hubiéramos ejecutado la app en local anteriormente y ya tuviésemos una carpeta \verb,node_modules,) \cite{stack-overflow:node_modules-folder}.
		\item \textbf{Se mantiene una entrada abierta para el contenedor}, lo que nos puede servir para poder instalar paquetes usando \verb,npm,.
\end{itemize}

\subsubsection{Backend}

\begin{lstlisting}
  api-rest:
    build:
      context: ../webapps/apirest
    restart: always
    ports:
      - ${API_WEB_DOCKER_PORT}:${API_WEB_LOCAL_PORT}
      - ${API_RELOAD_DOCKER_PORT}:${API_RELOAD_LOCAL_PORT}
      - ${API_DEBUG_DOCKER_PORT}:${API_DEBUG_LOCAL_PORT}
    networks:
      - nutrimenu-net
    env_file:
      - .env
    depends_on:
      mysqldb:
        condition: service_healthy
    working_dir: /app
    command: sh run_docker_dev.sh
    volumes:
      - ../webapps/apirest:/app
      - .m2:/root/.m2
\end{lstlisting}

Para la parte del \textit{backend} muchos de los elementos del Dockerfile son similares a los vistos en la parte del \textit{frontend}, así que simplemente destacaré aquellos que varían respecto a este:

\begin{itemize}
	\item Se van a \textbf{configurar los puertos} descritos en el apartado \ref{entorno-desarrollo}.
	\item A la hora de depender del contenedor de la base de datos para arrancar, he especificado la condición \verb,service_healthy,, que hace que no sólo el contenedor de la API espere a que el de la base de datos esté en estado \textit{ready}, sino que además \textbf{va a ejecutar un \textit{healthcheck} para comprobar que la base de datos se encuentra accesible}. Esto es \textbf{MUY importante} para el correcto funcionamiento de la app, ya que sin este \textit{healthcheck}, se puede dar el caso de que la API REST arranque antes que la base de datos, causando un error en la ejecución.
	\item Nada más arrancar \textbf{va a ejecutar el comando} \verb,run_docker_dev.sh,. Este comando lo que hace es levantar el servicio de \textit{backend} y estar a su vez a la espera de cambios en los ficheros del código fuente para, en caso de que suceda, recompilar el código de forma transparente. Este apartado es distinto en el archivo \verb,docker-compose.ci.yml,, puesto que hace referencia a \verb,run_docker_ci.sh, que simplemente compila la aplicación y la ejecuta.
	\item Se van a conectar dos volúmenes al contenedor: uno va a ser \textbf{el que contiene el código del servicio} (ya que si simplemente copiáramos el código al contenedor, cada vez que hubiera cambios en el código habría que volver a desplegar los contenedores, y de esta forma podemos realizar cambios en tiempo real); y otro va a ser \textbf{el volumen encargado de almacenar las dependencias de Maven}, para no tener que descargarlas cada vez que ejecutemos el código.
\end{itemize}

\subsubsection{Base de datos}

\begin{lstlisting}
  mysqldb:
    image: "mysql:8.0"
    restart: always
    ports:
      - ${DB_DOCKER_PORT}:${DB_LOCAL_PORT}
    networks:
      - nutrimenu-net
    env_file:
      - .env
    volumes:
      - ../data:/docker-entrypoint-initdb.d
      - db_data:/var/lib/mysql
    healthcheck:
      test: ["CMD", "mysqladmin" ,"ping", "-h", "localhost", "-u", "root", "-p${MYSQL_ROOT_PASSWORD}"]
      interval: 15s
      timeout: 5s
      start_period: 10s
      retries: 3
\end{lstlisting}

La base de datos también va a disponer de su propio contenedor, con ciertos parámetros similares a los de los servicios del \textit{backend}. Podemos destacar de aquí algunos elementos:

\begin{itemize}
	\item Vamos a usar la \textbf{imagen oficial de MySQL} (\href{https://hub.docker.com/layers/library/mysql/8/images/sha256-1717a37c2adc7772b569eb727bb198b6792f7d50c2adb4839ba637b5c9d88f42?context=explore}{Docker Hub}), concretamente la versión \textbf{8} (con los últimos parches). 
	\item Se va a conectar un volumen a la base de datos, que no es más que \textbf{el volumen encargado de almacenar todos los datos} de la base de datos. Esto es para no tener que inicializar la base de datos cada vez que la despleguemos, y mantener así persistencia entre sesiones. Además, si no existen datos en la base de datos se ejecutará el script SQL localizado en \textit{data}, que contiene un conjunto de datos de prueba.
	\item Se ha definido el \textit{healthcheck} del que hemos hablado en el apartado del \textit{backend}, que lanza un comando para comprobar que el servicio esté totalmente arrancado. Este comando se va a lanzar un máximo de 3 veces en un intervalo de 15 segundos entre una y otra, siendo la primera 10 segundos después de que arranque el contenedor, y con un timeout de 5 segundos.
\end{itemize} 

\subsubsection{Volúmenes}

\begin{lstlisting}
volumes:
  db_data:
\end{lstlisting}

Es el volumen creado para la base de datos, ya que a diferencia de los volúmenes creados para el \textit{backend}, en los que realmente hemos conectado dos carpetas de nuestro equipo a los contenedores, este volumen se crea dentro de Docker.

\subsubsection{Redes}

\begin{lstlisting}
networks:
  nutrimenu-net:
\end{lstlisting}

Esta es la red creada para conectar todos los contenedores y que no existan problemas de conectividad entre ellos.

\subsection{Configuración del debugger remoto}

Como veremos en el apartado \ref{entorno-desarrollo}, se van a configurar distintos puertos para el servicio de \textit{backend}, y uno de estos va a ser el puerto de \textit{debug}. Como vamos a ejecutar la aplicación usando Docker Compose y no compilando el código directamente desde nuestro IDE, por defecto no podemos hacer debug como haríamos con cualquier código compilado localmente. 

Sin embargo, con los pasos que vamos a ver a continuación \cite{medium:spring-boot-in-docker}, \textbf{vamos a ser capaces de conectarnos a nuestro entorno \textit{contenerizado} y realizar debug del código como lo haríamos de forma normal}. En este caso voy a demostrar cómo funcionaría en IntelliJ IDEA, pero mientras que el IDE disponga de un \textit{debugger remoto} y sigamos correctamente las instrucciones del desarrollador del IDE para usarlo, no debería de haber ningún problema.

\imagen{IntelliJ/debug-edit-conf-1}{Listado de todas las configuraciones de ejecución disponibles}{0.8}

Para comenzar, debemos de abrir  el proyecto en IntelliJ y en la parte superior, al lado de los botones de Ejecución y Debug, tenemos un selector donde podemos seleccionar la configuración que deseemos. Debemos hacer click en \textit{Edit Configurations...}.

\imagen{IntelliJ/debug-edit-conf-2}{Selección del tipo de configuración}{0.8}

Debemos hacer click en el icono con un + para añadir una nueva configuración, y seleccionar dentro del desplegable \textit{Remote JVM Debug}.

\imagen{IntelliJ/debug-edit-conf-3}{Creación de la configuración del debugger remoto}{0.8}

A continuación daremos nombre a la configuración, indicaremos el puerto correspondiente al servicio que estemos configurando, y seleccionaremos el módulo de dicho servicio.

\imagen{IntelliJ/debug-edit-conf-4}{Creación de un módulo para un servicio}{0.6}

En caso de que el módulo correspondiente no aparezca, deberemos seleccionar en la barra de herramientas \textit{File > New > Module from Existing Sources...}, seleccionar el directorio \verb,/webapps/apirest,, y escoger Maven dentro de \textit{Import module from external model}. Si sí que nos aparece el módulo podemos ignorar este paso.

\imagen{IntelliJ/debug-edit-conf-5}{Debug del servicio}{1}

Una vez hayamos completado estos pasos, podemos guardar la configuración y ejecutar el proyecto (asegurándonos de que la configuración deseada está seleccionada) desde el botón de Debug, como haríamos de normal. Si todo ha ido bien, el IDE nos devolverá en la consola el mensaje \textit{Connected to the target VM, address: 'localhost:5005', transport: 'socket'}.

\section{Compilación, despliegue y ejecución del proyecto}

Para realizar la compilación, despliegue y ejecución del proyecto es imprescindible disponer de \textbf{Docker CLI} y \textbf{Docker Compose} instalados en el equipo. La mejor forma de instalarlos es mediante \textbf{Docker Desktop}, ya que se encarga de instalar en el equipo el \textit{daemon} de Docker, Docker CLI, Docker Compose, y demás dependencias y herramientas que no vamos a utilizar para este proyecto, pero que nos pueden ayudar en el futuro. Docker Desktop es compatible con Windows, macOS y Linux, y soporta tanto arquitecturas x86 como ARM64 (Apple Silicon).

\subsection{Instalación de Docker Desktop}

\imagen{Docker/docker-download}{Página de descarga de Docker Desktop}{0.8}

Para realizar la instalación de Docker Desktop tan sólo debemos dirigirnos a \url{https://docs.docker.com/get-docker/}, seleccionar la plataforma deseada, y descargar el instalador.

\imagen{Docker/docker-installation-mac}{Instalación de Docker Desktop en Mac}{0.8}

En caso de realizar la instalación en Windows, es importante tener en cuenta que se debe tener WSL (\textit{Windows Subsystem for Linux}) instalado y con una distribución activa, ya que si no la instalación puede dar errores (figura \ref{fig:Docker/docker_error_wsl}). Para realizar esto se puede usar \href{https://learn.microsoft.com/es-es/windows/wsl/install}{la guía de Microsoft}. También se puede seguir \href{https://learn.microsoft.com/es-es/windows/wsl/tutorials/wsl-containers}{esta guía} para usar Docker desde WSL.

\imagen{Docker/docker_error_wsl}{Error de instalación al no tener WSL instalado}{0.8}

Una vez ejecutemos el instalador (en caso de macOS simplemente se debe arrastrar la aplicación a la carpeta de Aplicaciones), y hayamos seguido todos los pasos hasta finalizar la instalación, nos encontraremos con el panel principal de Docker Desktop. 

\imagen{Docker/docker-main-window}{Ventana principal de Docker Desktop}{0.8}

En esta ventana podemos ver todos los contenedores que están ejecutándose actualmente en el sistema, así como las imágenes descargadas y los volúmenes creados.

\subsection{Preparación del entorno de desarrollo} \label{entorno-desarrollo}

Para hacer más sencillo el cambio de valores en distintos parámetros de la aplicación, como el mapeado de puertos o la gestión de secretos, y facilitar así el despliegue de nuevos entornos o la escalabilidad de los servicios, estos valores se han externalizado en un sólo archivo \verb,.env,, que va a ser un archivo que simplemente va a contener las variables junto a sus valores. 

Al ser un fichero que contiene secretos, siguiendo lo que dictan las buenas prácticas no se sincronizará con el repositorio, ya que si se tratase de un entorno de producción o el código se hiciera público, cualquier usuario podría obtener acceso a los datos de los usuarios de la aplicación.

Por lo tanto, antes de desplegar los contenedores, para que la aplicación funcione correctamente \textbf{se debe crear este archivo} \verb,.env, \textbf{dentro del directorio} \verb,/docker-compose,, y se deben de definir las siguientes variables dentro de él:

\begin{lstlisting}[language=Bash]
# Variables de la API REST
API_WEB_DOCKER_PORT=8080
API_WEB_LOCAL_PORT=8080
API_RELOAD_DOCKER_PORT=35729
API_RELOAD_LOCAL_PORT=35729
API_DEBUG_DOCKER_PORT=5005
API_DEBUG_LOCAL_PORT=5005

# Variables del frontend
CHOKIDAR_USEPOLLING=true
VITE_NUTRI_X_APP_ID=<YOUR_APP_ID> # Introduce el App ID de la API de Nutritionix
VITE_NUTRI_X_APP_KEY=<YOUR_APP_KEY> # Introduce el App Key de la API de Nutritionix
VITE_NUTRI_X_REMOTE_USER_ID=0

# Variables de la base de datos
DATABASE_HOST=mysqldb
MYSQL_DATABASE=nutri_db
MYSQL_USER=dbadmin
MYSQL_PASSWORD=<MySQL password> # Introduce la password que desees para el usuario dbadmin
MYSQL_ROOT_PASSWORD=<MySQL root password> # Introduce la password que desees para el usuario root
DB_DOCKER_PORT=3306
DB_LOCAL_PORT=3306
\end{lstlisting}

\textbf{NOTA:} Para crear un documento oculto en Windows, como un fichero \verb,.env,, podemos usar \href{https://notepad-plus-plus.org/}{Notepad++}. Para ello, abriremos el programa, escribiremos el contenido del fichero \verb,.env,, daremos a \verb,Save > Save As,, seleccionaremos en tipo de fichero \verb,All files,, y nombraremos el fichero como \verb,.env, \cite{stack-overflow:env-file-windows}.

Los valores pueden ser modificados si es necesario, pero se recomienda usar los valores dados (\textbf{es necesario dar los valores que queramos en las variables} \verb,MYSQL_PASSWORD, y \verb,MYSQL_ROOT_PASSWORD,). Con estos valores lo que estamos indicando es:

\begin{itemize}
	\item Los puertos en los que se van a ejecutar cada uno de los servicios dentro de los contenedores (\verb,LOCAL_PORT,) y el puerto desde el cuál vamos a acceder al servicio correspondiente (\verb,DOCKER_PORT,), que en este caso no hay necesidad de diferenciarlos, y en este ejemplo son el \textbf{8080} para la API REST, el \textbf{3000} para el frontend y el \textbf{3306} para la base de datos.
	\item El puerto en el que va a estar escuchando el plugin \textit{LiveReload} de Spring, que es el encargado de recargar el servidor en caliente cada vez que se produzcan cambios en el código, y no tener así que recompilar todo el código cada vez que cambiemos algo. En este caso el puerto va a ser el \textbf{35729}.
	\item El puerto que se va a usar para permitir conectar un debugger remoto a la aplicación, y poder así hacer debug desde el IDE que usemos para el desarrollo. El puerto va a ser el \textbf{5005} para la API REST.
	\item Configuración adicional de la parte del frontend, como la variable \verb,CHOKIDAR_USEPOLLING, que indica que usemos el mecanismo \textit{CHOKIDAR} para observar cambios en los archivos fuente, para así poder aplicar los cambios de forma automática y realizar una recarga de la página de forma instantánea. Es necesario usar este mecanismo en aplicaciones React que se encuentren desplegadas en VMs o dockerizadas, como es este caso \cite{create-react-app:advanced-conf}.
	\item \textbf{Las claves para poder conectar con la API de Nutritionix}. Esto es esencial para que la búsqueda de alimentos funcione, ya que sin esto sólo se pueden crear alimentos a mano.
	\item Todas las variables correspondientes a la base de datos, como el \textbf{nombre de la instancia}, \textbf{nombre de la base de datos}, y \textbf{credenciales de los usuarios de MySQL}.
\end{itemize}

% TODO añadir cómo obtener credenciales de Nutritionix

\subsection{Despliegue del entorno de desarrollo}

\imagen{Docker/docker-compose-build}{Creación de las imágenes de los contenedores Docker}{0.8}

Una vez definidas las variables de entorno, podemos probar que todo funciona correctamente construyendo las imágenes de los contenedores de los servicios. Para ello, abriremos una consola (se puede usar PowerShell en Windows, o la nueva Terminal de Windows 11) y accederemos al directorio \verb,/docker-compose,, (\textbf{TODOS los comandos que veremos a continuación deben ejecutarse dentro del directorio} \verb,docker-compose,), donde ejecutaremos el siguiente comando:

\begin{lstlisting}[language=Bash]
docker compose -f docker-compose.dev.yml build
\end{lstlisting}

\imagen{Docker/docker-compose-up}{Despliegue de los contenedores Docker}{0.8}

Este proceso de construcción de imágenes sólo es necesario realizarlo la primera vez que ejecutemos el proyecto, o en caso de realizar cambios en los Dockerfiles, ya que estaríamos cambiando la estructura de las imágenes. Como tanto la capa de \textit{frontend} como la de \textit{backend} disponen de mecanismos para recargarse automáticamente en caso de cambios en el código, no hace falta reiniciar los contenedores cada vez que realicemos cambios en las aplicaciones.

Una vez haya terminado de construir las imágenes, para desplegar los contenedores y levantar el servicio ejecutaremos el siguiente comando:

\begin{lstlisting}[language=Bash]
docker compose -f docker-compose.dev.yml up -d
\end{lstlisting}

\imagen{Docker/docker-desktop-containers}{Ejecución multi-contenedor de la aplicación}{0.8}

Si todo ha ido bien, podremos ver en Docker Desktop los contenedores corriendo con estado \textit{Running}, y ya podremos acceder a la aplicación como haríamos con cualquier tipo de despliegue. 

\imagen{Docker/docker-desktop-logs}{Logs de la ejecución de un contenedor en Docker Desktop}{0.8}

Para asegurarnos de que no hay ningún tipo de problema, podemos seleccionar los 3 puntos verticales situados a la derecha de cada contenedor, y hacer click en \textit{View details}, ya que esto nos permite poder ver el log de cada servicio en tiempo real. 

\imagen{Docker/docker-compose-stop}{Parada de los contenedores en ejecución}{0.8}

En caso de querer parar la ejecución de los contenedores, podemos hacerlo con el siguiente comando:

\begin{lstlisting}[language=Bash]
  docker compose -f docker-compose.dev.yml stop
\end{lstlisting}

\imagen{Docker/docker-desktop-stop}{Contenedores parados pero no borrados en Docker Desktop}{0.8}

Este comando lo que va a hacer es parar los servicios, pero no borrarlos, pudiendo reanudar el servicio otra vez en cualquier momento con \verb,docker compose -f docker-compose.dev.yml up -d,.

\imagen{Docker/docker-compose-down}{Parada y borrado de los contenedores en ejecución}{0.8}

En caso de querer borrar los contenedores, en vez de usar el comando que hemos usado arriba para parar los servicios, usaríamos el siguiente comando, que para y elimina los contenedores:

\begin{lstlisting}[language=Bash]
  docker compose -f docker-compose.dev.yml down
\end{lstlisting}

\imagen{Docker/docker-desktop-down}{Contenedores de Docker totalmente eliminados}{0.8}

\imagen{Docker/docker-desktop-down-2}{Persistencia del volumen de la base de datos}{0.8}

En el caso de nuestra aplicación, como el código se monta desde nuestro equipo en todos los servicios como un volumen externo, y la base de datos se encuentra en un volumen de Docker, borrar los contenedores no supondría ninguna pérdida de datos, así que realmente no habría diferencia respecto a simplemente parar los servicios, al menos a nivel de datos. Si reconstruyésemos las imágenes entonces sí que tendría más sentido borrar los contenedores primero, ya que así la próxima vez que los desplegásemos lo harían con la nueva imagen.

\imagen{Docker/docker-compose-down-volumes}{Parada y borrado de contenedores y volúmenes existentes}{0.8}

En caso de querer borrar los contenedores y los volúmenes asociados (en nuestro caso, el volumen de la base de datos), deberíamos de ejecutar el siguiente comando:

\begin{lstlisting}[language=Bash]
  docker compose -f docker-compose.dev.yml down --volumes
\end{lstlisting}

\imagen{Docker/docker-desktop-down-volumes}{Volumen de la base de datos totalmente eliminado}{0.8}

Hay que tener en mente que ejecutar este comando \textbf{hará que se pierdan todos los datos de la base de datos}.

\subsection{Acceso a la aplicación}

Si se han mantenido los puertos recomendados en el archivo .env de ejemplo, podremos acceder a la aplicación desde el navegador escribiendo \textbf{http://localhost:3000/}. Para poder interactuar con la aplicación se han creado tres usuarios distintos, cada uno con un rol distinto, cuyas credenciales son:

\textbf{Usuario de tipo administrador}
\begin{itemize}
	\item \textbf{Nombre de usuario}: administrador
	\item \textbf{Contraseña}: administrador1	
\end{itemize}

\textbf{Usuario de tipo editor}
\begin{itemize}
	\item \textbf{Nombre de usuario}: editorubu
	\item \textbf{Contraseña}: editorubu1	
\end{itemize}

\textbf{Usuario de tipo camarero}
\begin{itemize}
	\item \textbf{Nombre de usuario}: camareroubu
	\item \textbf{Contraseña}: camareroubu1	
\end{itemize}


\section{Pruebas del sistema}

A lo largo de la aplicación existen varias comprobaciones que permiten indicar al usuario que la acción que está realizando no está permitida, o que hay algo que no es correcto:

\begin{itemize}
	\item \textbf{Comprobación de que el usuario no existe}: Esta comprobación se puede encontrar tanto en el inicio de sesión como en el panel de contraseña olvidada.

\imagen{Testing/user_not_exists}{Mensaje de que no existe el usuario indicado}{0.8}

\item \textbf{Comprobación de que la contraseña es incorrecta}: Esta comprobación se puede encontrar al iniciar sesión.

\imagen{Testing/incorrect_password}{Mensaje de que la contraseña introducida es incorrecta}{0.8}
	
	\item \textbf{Comprobación de que las contraseñas no coinciden}: Esta comprobación se puede encontrar en el panel de contraseña olvidada, cuando se pide introducir la nueva contraseña dos veces.
	
\imagen{Testing/passwords_dont_match}{Mensaje de que las contraseñas no coinciden}{0.8}

\item \textbf{Comprobación de que un elemento requerido no se ha rellenado}: Esta comprobación se puede encontrar al crear elementos de distinto tipo, cuando no rellenamos un campo obligatorio.
	
\imagen{Testing/required_value}{Mensaje de que el valor es requerido}{0.8}

\item \textbf{Comprobación de que ya existe un elemento con ese valor}: Esta comprobación se puede encontrar al crear elementos de distinto tipo, cuando se intenta repetir un valor que es clave primaria de ese elemento.
	
\imagen{Testing/repeated_value}{Mensaje de que el valor se ha repetido}{0.8}

\item \textbf{Comprobación de que ya existe un alimento con ese nombre en la base de datos}: Esta comprobación se puede encontrar al buscar un alimento en Nutritionix.	
\imagen{Testing/alimento_exists}{Mensaje de que el alimento ya existe}{0.8}
	
\end{itemize}


\apendice{Documentación de usuario}

\section{Introducción}

Esta sección pretende ilustrar cómo funciona la aplicación, cuales son las distintas partes que la componen, y cómo el usuario puede interactuar con ellas. 

Para el desarrollo de este manual se parte con la idea de que la aplicación se encuentra desplegada y disponible.

\section{Requisitos de usuarios}

Como se trata de una aplicación web, los únicos requisitos por parte del usuario son:

\begin{itemize}
  \item Una conexión a Internet estable, que permita conectar con el servidor que almacena la aplicación (en caso de que el despliegue de la aplicación no sea local).
  \item Disponer de cualquier navegador web moderno, como Google Chrome, Mozilla Firefox, Safari o similares. No se recomienda usar un navegador desactualizado, ya que puede no ser compatible con algunas de las funcionalidades que ofrece React.
\end{itemize}


\section{Instalación}

La aplicación no requiere de ningún tipo de instalación, con disponer de un navegador web es suficiente.

\section{Manual del usuario}

Se van a proceder a detallar los distintos apartados que conforman la aplicación web:

\subsection{Vista inicial}

Nada más abrir la aplicación, lo primero que nos encontramos es la vista inicial. Desde aquí podemos usar la aplicación como un consumidor, o podemos pasar a la ventana de inicio de sesión haciendo click en el botón \textit{Iniciar sesión}.

\imagen{Manual/main}{Vista inicial de NutriMenu}{0.8}

\subsection{Inicio de sesión}

Esta es la vista desde la cual vamos a acceder a los paneles de gestión. Para acceder debemos de introducir las credenciales y hacer click en \textit{Inicio de sesión}, o hacer click en \textit{¿Has olvidado tu contraseña?} en caso de no recordarla.

\imagen{Manual/login}{Vista de inicio de sesión}{0.8}

\subsection{Recuperación de contraseña}

Para recuperar la contraseña se debe introducir el nombre de usuario de la cuenta que queremos recuperar, y a continuación se nos va a solicitar que escribamos la nueva contraseña dos veces, para asegurarnos de que la nueva contraseña es la que realmente queremos. Una vez la tengamos, debemos hacer click en Continuar, para que así esta contraseña se guarde y volvamos a la página de inicio de sesión.

\imagen{Manual/forgotpass_user}{Primera vista en el proceso de recuperar contraseña}{0.8}

\imagen{Manual/forgotpass_passwords}{Vista para establecer una nueva contraseña}{0.8}

\subsection{Panel de control de administradores}

Esto es lo primero que veremos al iniciar sesión con una cuenta de administrador. Podemos gestionar empresas, locales, usuarios y alimentos desde este tipo de cuenta.

\imagen{Manual/admin_panel}{Vista del panel de control de administradores}{0.8}

\subsection{Gestión de empresas}

En la sección de Gestión de empresas dispondremos de todas las empresas disponibles en la aplicación.

\imagen{Manual/empresa_table}{Vista principal de gestión de empresas}{0.8}

Para crear una nueva empresa, debemos hacer click en \textit{Crear una nueva empresa} dentro del panel de gestión de empresas. Debemos rellenar los datos que se marcan con asterisco, y cuando hayamos terminado, haremos click en \textit{Guardar}.

\imagen{Manual/empresa_new}{Vista de creación de empresas}{0.8}

Para editar una empresa, haremos click en \textit{Editar} en el panel de Gestión de empresas. Modificaremos los datos que deseemos y haremos click en \textit{Guardar} para que se apliquen los cambios.

\imagen{Manual/empresa_edit}{Vista de edición de empresas}{0.8}

Para borrar una empresa, haremos click en \textit{Borrar} en el panel de Gestión de empresas. Nos aparecerá una ventana pidiéndonos confirmar la acción, y si le damos a \textit{Sí, borrar} se producirá el borrado de la empresa.

\imagen{Manual/empresa_delete}{Vista de borrado de empresas}{0.8}

\subsection{Gestión de locales}

En la sección de Gestión de locales dispondremos de todos los locales disponibles en la aplicación.

\imagen{Manual/locales_table}{Vista principal de gestión de locales}{0.8}

Para crear un nuevo local, debemos hacer click en \textit{Crear un nuevo local} dentro del panel de gestión de locales. Debemos rellenar los datos que se marcan con asterisco, y cuando hayamos terminado, haremos click en \textit{Guardar}.

\imagen{Manual/locales_new}{Vista de creación de locales}{0.8}

Para editar un local, haremos click en \textit{Editar} en el panel de Gestión de locales. Modificaremos los datos que deseemos y haremos click en \textit{Guardar} para que se apliquen los cambios.

\imagen{Manual/locales_edit}{Vista de edición de locales}{0.8}

Para borrar un local, haremos click en \textit{Borrar} en el panel de Gestión de locales. Nos aparecerá una ventana pidiéndonos confirmar la acción, y si le damos a \textit{Sí, borrar} se producirá el borrado del local.

\imagen{Manual/locales_delete}{Vista de borrado de locales}{0.8}

\subsection{Gestión de usuarios}

En la sección de Gestión de usuarios dispondremos de todos los usuarios disponibles en la aplicación.

\imagen{Manual/usuarios_table}{Vista principal de gestión de usuarios}{0.8}

Para crear un nuevo usuario, debemos hacer click en \textit{Crear un nuevo usuario} dentro del panel de gestión de usuarios. Debemos rellenar los datos que se marcan con asterisco, y cuando hayamos terminado, haremos click en \textit{Guardar}. En caso de que el usuario sea de tipo Administrador, no hace falta asignar una empresa al usuario, pero si tiene un rol distinto sí se debe hacer.

\imagen{Manual/usuarios_new}{Vista de creación de usuarios}{0.8}

Para editar un usuario, haremos click en \textit{Editar} en el panel de Gestión de usuarios. Modificaremos los datos que deseemos y haremos click en \textit{Guardar} para que se apliquen los cambios.

\imagen{Manual/usuarios_edit}{Vista de edición de usuarios}{0.8}

Para borrar un usuario, haremos click en \textit{Borrar} en el panel de Gestión de usuarios. Nos aparecerá una ventana pidiéndonos confirmar la acción, y si le damos a \textit{Sí, borrar} se producirá el borrado del usuario.

\imagen{Manual/usuarios_delete}{Vista de borrado de usuarios}{0.8}

\subsection{Gestión de alimentos}

En la sección de Gestión de alimentos dispondremos de todos los alimentos disponibles en la aplicación.

\imagen{Manual/alimentos_table}{Vista principal de gestión de alimentos}{0.8}

Para buscar un alimento en la base de datos de Nutritionix, debemos hacer click en \textit{Crear un nuevo alimento} dentro del panel de gestión de usuarios. A continuación, debemos asegurarnos que estamos en la sección \textit{Buscar en Nutritionix}, aunque es la sección que se abre por defecto. 

Introduciremos el nombre del alimento a buscar, y haremos click en \textit{Buscar}. El nombre del alimento no tiene por qué ser exacto, ya que la búsqueda soporte lenguaje natural, aunque cuanto más precisos seamos, mejor serán los resultado.

Si existen alimentos que respondan a la búsqueda que acabamos de hacer, se mostrarán en forma de lista junto a una imagen de este mismo. Si queremos añadir alguno de estos alimentos haremos click en \textit{Añadir}, y el alimento se añadirá a nuestra base de datos. La aplicación nos redirigirá al panel de Gestión de alimentos al terminar.

\imagen{Manual/alimentos_nutritionix}{Vista de búsqueda de alimentos en Nutritionix}{0.8}

Si por el contrario queremos crear un alimento a mano, deberemos de hacer click en \textit{Crear a mano}. Primero se nos pedirá la información básica del alimento, así que deberemos de rellenar los datos pedidos y hacer click en \textit{Continuar}. 

A continuación veremos varias páginas de introducción de datos, las correspondientes a la introducción de datos de componentes, a introducción de vitaminas, y a introducción de minerales. Una vez estén todas rellenas y hagamos click en \textit{Guardar}, el alimento se guardará en la base de datos y la aplicación nos devolverá al panel de Gestión de elementos.

\imagen{Manual/alimentos_create}{Vista de creación de alimentos a mano}{0.8}

\imagen{Manual/alimentos_nutrients}{Vista para indicar los componentes de un alimento}{0.8}

\imagen{Manual/alimentos_vitamins}{Vista para indicar las vitaminas de un alimento}{0.8}

\imagen{Manual/alimentos_minerals}{Vista para indicar los minerales de un alimento}{0.8}

Si dentro del panel de Gestión de alimentos hacemos click en el nombre de un alimento, se nos dirigirá a la vista de componentes nutricionales de un alimento. Aquí podemos observar tanto el gráfico de componentes como la tabla detallada de componentes nutricionales del alimento.

\imagen{Manual/alimentos_view_1}{Vista de componentes nutricionales de un alimento}{0.8}

En caso de poner el ratón encima de alguno de los quesos del gráfico, podemos ver los valores exactos de ese componente.

\imagen{Manual/alimentos_view_2}{Vista con la distinta información que muestran las tablas de componentes}{0.8}

\imagen{Manual/alimentos_view_details}{Vista con información adicional en el gráfico}{0.8}

Para editar un alimento, haremos click en \textit{Editar} en el panel de Gestión de alimentos. Modificaremos los datos que deseemos y haremos click en \textit{Guardar} para que se apliquen los cambios.

\imagen{Manual/alimentos_edit_1}{Vista de edición de alimentos}{0.8}

\imagen{Manual/alimentos_edit_2}{Vista de edición de los componentes de un alimento}{0.8}

\imagen{Manual/alimentos_edit_3}{Vista de edición de las vitaminas de un alimento}{0.8}

\imagen{Manual/alimentos_edit_4}{Vista de edición de los minerales de un alimento}{0.8}

Para borrar un alimento, haremos click en \textit{Borrar} en el panel de Gestión de alimentos. Nos aparecerá una ventana pidiéndonos confirmar la acción, y si le damos a \textit{Sí, borrar} se producirá el borrado del alimento.

\imagen{Manual/alimentos_delete}{Vista de borrado de alimentos}{0.8}

\subsection{Panel de control de editores}

Esto es lo primero que veremos al iniciar sesión con una cuenta de editor. Podemos gestionar los menús y platos de la empresa a la que pertenece el usuario desde este tipo de cuenta.

\imagen{Manual/editor_panel}{Vista del panel de control de editores}{0.8}

\subsection{Gestión de platos}

En la sección de Gestión de locales dispondremos de todos los locales disponibles en la aplicación.

\imagen{Manual/platos_table}{Vista principal de gestión de platos}{0.8}

Para crear un nuevo plato, debemos hacer click en \textit{Crear un nuevo plato} dentro del panel de gestión de platos. Debemos rellenar la información básica del plato en el panel \textit{Información sobre el plato}, y a continuación podemos proceder a añadir los alimentos que van a formar el plato.

Para añadir alimentos a un plato existen tres formas:

\begin{itemize}
	\item Añadir el alimento desde la lista de alimentos existentes en la base de datos
	\item Buscar un alimento en Nutritionix
	\item Crear un alimento a mano	
\end{itemize}

Según vayamos añadiendo los alimentos al plato, estos irán apareciendo en el panel \textit{Lista de alimentos}. Podemos modificar su cantidad en gramos, o borrarlos del plato.

Cuando estemos satisfechos con el plato, debemos hacer click en \textit{Guardar}. En caso de querer cancelar la creación del plato, debemos hacer click en \textit{Cancelar}, donde se nos abrirá una ventana preguntando si realmente queremos cancelar la creación del plato.

\imagen{Manual/platos_new_table}{Vista de creación de platos}{0.8}

\imagen{Manual/platos_new_nutritionix}{Vista para añadir un alimento de Nutritionix al plato}{0.8}

\imagen{Manual/platos_new_create}{Vista para crear un alimento y añadirlo al plato}{0.8}

\imagen{Manual/platos_new_cancel}{Vista de cancelación de la creación de un plato}{0.8}

Para editar un plato, haremos click en \textit{Editar} en el panel de Gestión de platos. Modificaremos los datos que deseemos y haremos click en \textit{Guardar} para que se apliquen los cambios.

\imagen{Manual/platos_edit}{Vista de edición de platos}{0.8}

Para borrar un plato, haremos click en \textit{Borrar} en el panel de Gestión de platos. Nos aparecerá una ventana pidiéndonos confirmar la acción, y si le damos a \textit{Sí, borrar} se producirá el borrado del plato.

\imagen{Manual/platos_delete}{Vista de borrado de platos}{0.8}

Si dentro del panel de Gestión de platos hacemos click en el nombre de un plato, se nos dirigirá a la vista de componentes nutricionales de un plato. Aquí podemos observar tanto el gráfico de componentes como la tabla detallada de componentes nutricionales del plato.

Debajo del informe nutricional también disponemos de los alimentos que componen el plato, así como su cantidad. Si hacemos click en el nombre de un alimento podemos ver el informe nutricional de ese alimento.

\imagen{Manual/platos_view}{Vista de componentes nutricionales e ingredientes de un plato}{0.8}

\subsection{Gestión de menús}

En la sección de Gestión de menús dispondremos de todos los menús disponibles en la aplicación.

\imagen{Manual/menus_table}{Vista principal de gestión de menús}{0.8}

Para crear un nuevo menú, debemos hacer click en \textit{Crear un nuevo menú} dentro del panel de gestión de menús. Debemos rellenar la información básica del menú en el panel \textit{Información sobre el menú}, y a continuación podemos proceder a seleccionar los locales en los que va a aparecer este menú y los platos que queremos añadir en él.

Cuando estemos satisfechos con el menú, debemos hacer click en \textit{Guardar}. En caso de querer cancelar la creación del menú, debemos hacer click en \textit{Cancelar}, donde se nos abrirá una ventana preguntando si realmente queremos cancelar la creación del plato.

\imagen{Manual/menus_new}{Vista de creación de menús}{0.8}

\imagen{Manual/menus_new_cancel}{Vista de cancelación de la creación de un menú}{0.8}

Para editar un menú, haremos click en \textit{Editar} en el panel de Gestión de menús. Modificaremos los datos que deseemos y haremos click en \textit{Guardar} para que se apliquen los cambios.

\imagen{Manual/menus_edit}{Vista de edición de menús}{0.8}

Para borrar un menú, haremos click en \textit{Borrar} en el panel de Gestión de menús. Nos aparecerá una ventana pidiéndonos confirmar la acción, y si le damos a \textit{Sí, borrar} se producirá el borrado del menú.

\imagen{Manual/menus_delete}{Vista de borrado de menús}{0.8}

Si dentro del panel de Gestión de menús hacemos click en el nombre de un menú, se nos dirigirá a la vista de componentes nutricionales de un menú. Esta es la misma vista que se van a encontrar los consumidores al seleccionar un menú.

Si seleccionamos un plato, podemos observar como automáticamente se actualizan tanto el gráfico como la tabla (podemos alternar entre ellos pulsando en los botones \textit{Gráfico} y \textit{Tabla de nutrientes}). Podemos seleccionar o deseleccionar tantos platos como deseemos, ya que el informe se va a actualizar de forma automática.

\imagen{Manual/menus_view_1}{Vista de informe nutricional de un menú (gráfica)}{0.8}

En caso de que un valor supere el 80\% de su cantidad recomendada, la fila de ese valor se pondrá de color amarillo. En caso de que ese valor supere su cantidad recomendad, la fila se pondrá de color rojo.

\imagen{Manual/menus_view_2}{Vista de informe nutricional de un menú (tabla de nutrientes)}{0.8}

\subsection{Panel de control de camareros}

Esto es lo primero que veremos al iniciar sesión con una cuenta de camarero. Podemos gestionar los menús de la empresa a la que pertenece el usuario desde este tipo de cuenta.

\imagen{Manual/camarero_view}{Vista del panel de control de camareros}{0.8}

\subsection{Parte para consumidores}

Esta es la parte de la aplicación a la que van a acceder los consumidores. Lo primero que nos pregunta la aplicación es que seleccionemos una empresa.

\imagen{Manual/main}{Vista de selección de una empresa}{0.8}

A continuación se nos pide que escojamos un local de los pertenecientes a la empresa que acabamos de escoger.

\imagen{Manual/seleccion_local}{Vista de selección de un local}{0.8}

Por último, se nos pide que seleccionemos una fecha para ver los menús disponibles. Por defecto la fecha escogida es la del día de hoy, pero se puede modificar para escoger otra fecha.

Una vez hayamos seleccionado la fecha deseada, podemos ver los menús publicados para ese día. Si hacemos click en alguno de ellos podemos ver el informe nutricional de ese menú, el cual funciona de la misma que acabamos de explicar en el panel de Gestión de menús.

\imagen{Manual/seleccion_menu}{Vista de selección de un menú}{0.8}

\imagen{Manual/menus_view_1}{Vista de selección de platos, junto a la gráfica nutricional}{0.8}

\imagen{Manual/menus_view_2}{Vista de selección de platos, junto a la tabla de nutrientes}{0.8}


\bibliographystyle{plain}
\bibliography{bibliografiaAnexos}

\end{document}
